\label{ch:introduction}

\section{Motivation}

Core diagnostics have been a key component of nuclear reactor operation to ensure reactor safety and performance. As a multicomponent system, perturbations in nuclear reactors are unavoidable. For power reactors, neutron noise is defined as the stochastic or random process that always happens in a nuclear reactor \cite{saitoTheoryPowerReactor1974b}. In the past, neutron noise experiments were developed to perform online monitoring and diagnostics of research reactors. Some examples include continuous neutron noise monitoring at the \gls*{HFIR}, neutron and pressure noise monitors at the \gls*{MSRE} \cite{fryExperienceReactorMalfunction1971}, neutron noise diagnostics at the Palisades Nuclear Generating Station in Michigan \cite{fryAnalysisNeutrondensityOscillations1975}. German measurements in \gls*{BWR} that showed vibrations quantified using neutron noise simulations \cite{wachInvestigationJointEffect1974}.

The success of core diagnostics using neutron noise experiments motivates the development of computational models of neutron noise. In the early days, the computational models revolved around accurately modeling perturbations in a reactor as cross-section perturbations \cite{saitoTheoryPowerReactor1974,saitoTheoryPowerReactor1974a}. Further development of the model led to the concept of the noise unfolding method, a technique for detecting the location of neutron noise and determining its magnitude. Three main methods have been developed to unfold neutron noise: the inversion method, the zoning method, and the scanning method. All of the methods require the Green function matrix to solve the problem \cite{pazsitNoiseTechniquesNuclear2010}. The input for these methods is the neutron noise derived from the detector readings, and the output is the noise locations and their magnitudes.

The development of advanced reactors has been rising for the last several years. This is supported by the backing of the financial market through startups and private funds, as well as a favorable public perception of nuclear energy in recent years. While many developers focus on the inherent reactor safety of advanced reactors, it is also essential to develop monitoring and diagnostic methods in advanced reactors.

In this dissertation, computational models of neutron noise are developed based on the neutron diffusion equation in the frequency domain. The solver is developed using the finite volume method for rectangular and triangular geometries. The main goal of the solver is application to \glspl*{HTGR}. However, application to rectangular geometries is also provided to differentiate the characteristics of \glspl*{HTGR} and \glspl*{LWR} systems. In this dissertation, code-to-code comparisons are provided. Methods for unfolding neutron noise are also developed in the simulator to highlight the advantages and disadvantages of the methods. 

\section{Objectives}

The main objective of this work is to develop and demonstrate core diagnostics in \gls*{HTGR} system using the neutron noise method. The following objectives are identified to advance this goal:
\begin{enumerate}
    \item Develop a numerical solver to perform forward, adjoint, and noise (in frequency domain) calculations, based on multigroup diffusion theory, in cartesian and triangular geometry.
    \item Determine the characteristics of noise in \glspl*{HTGR} compared to \glspl*{LWR} through the zero-power reactor transfer function.
    \item Determine the distinct characteristics of noise sources in \glspl*{LWR} and \glspl*{HTGR} for absorber of variable strength and fuel assembly vibration.
    \item Perform noise unfolding in prismatic \gls*{HTGR} using existing methods.
    \item Derive new noise unfolding methods and demonstrate their performance relative to the existing methods.
\end{enumerate}

Objective 1 focuses on the development of a numerical solver that could perform simulations in rectangular and hexagonal geometries using the provided macroscopic cross-sections. This solver will be used to model and simulate neutron noise in \gls*{HTGR} systems, which is the focus of Objective 2. The \gls*{HTGR} model used in this case is the \gls*{HTTR} benchmark model. The \gls*{HTTR} benchmark is modeled in the Monte Carlo code Serpent \cite{leppanenSerpentMonteCarlo2015}. This Serpent model is a heterogeneous model with various materials. The details of the \gls*{HTTR} model will be provided in the following sections. This model is mainly used to generate homogenized cross-sections for the solver. 

After obtaining the model, Objective 3 and 4 focus on modeling various neutron noise models and unfolding methods that existed in the literature. These models and methods are mainly developed for \gls*{LWR} systems. However, it is meaningful to apply the methods and model to understand the different behavior of \gls*{HTGR} system under small perturbations. The novelty of this work relies on Objective 5. Using all the information known from previous Objectives, it is important to extend these models and methods to reflect the important physics of \gls*{HTGR} system and leverage all the known parameters from forward and adjoint calculations.

Objectives 1, 2, and 3 are preliminary and background work necessary to advance beyond the current state-of-the-art. Objective 4 establishes baseline performance of existing noise unfolding methods, which is improved by new methods derived in Objective 5. Objectives 3, 4, and 5 all contribute a new body of knowledge. Objective 3 contributes a quantitative explanation of \glspl*{LWR} and \glspl*{HTGR} neutron physics differences. Objective 4 contributes a consistent performance comparison of existing noise unfolding methods for the \gls*{HTGR} reactor. Objective 5 derives new noise unfolding methods and demonstrates their superior performance relative to existing methods. Objective 5 is a new and original contribution to the body of knowledge.

\section{Relevant Works}

Relevant works and toolsets that are related to the objectives are presented in this section. These works are compared with the current work to identify the gaps in the method. The discussion is divided into two criteria. Namely, the neutron noise solver and the noise unfolding method.

For the neutron noise solver, this work focuses on solving the neutron noise equation using a deterministic approach, namely the diffusion method. This work primarily follows the CORE SIM+ simulator developed at Chalmers University of Technology, Sweden. CORE SIM+ is a 3D neutron diffusion solver that can solve forward, adjoint, and noise problems for rectangular geometries \cite{mylonakisCORESIMFlexible2021}. This simulator has been used to simulate neutron noise for experiments and has been validated using experimental data \cite{mylonakisCORESIMSimulations2021,hursinModelingNoiseExperiments2023}. CORE SIM+ uses a box-scheme finite difference for spatial discretization and frequency domain analysis. The difference between the simulator and this work is the energy treatment. This work uses a multigroup approach, while CORE SIM+ only uses two-group approach. Having a multigroup approach is generally beneficial to advanced reactors, especially \glspl*{HTGR}, since the neutron spectrum is different from \glspl*{PWR} \cite{ardiansyahEvaluationPBMR400Core2021}. Another code to be used for comparison is FEMFFUSION, which was developed in Spain as the time-domain counterpart of CORE SIM+ in the \gls*{CORTEX} project \cite{demaziereCORTEXProjectImproving2020}. FEMFFUSION uses time-domain for noise analysis and the finite element method for discretization. FEMFFUSION can also solve the neutron transport equation using simplified spherical harmonic ($\text{SP}_\text{N}$) approximations. Table \ref{table:other_works} summarizes the different approaches in the simulators.

\begin{table}[ht]
    \centering
    \caption{Summary of features of simulators for noise calculation}
    {\renewcommand{\arraystretch}{1.5}%
    \begin{tabular}{ | >{\raggedright\arraybackslash}m{3.0cm} | >{\centering\arraybackslash}m{3.5cm} |>{\centering\arraybackslash}m{3.0cm} | >{\centering\arraybackslash}m{3.5cm} | } 
     \hline
     \textbf{Parameter} & \textbf{CORE SIM+} & \textbf{FEMFFUSION}& \textbf{Current Work} \\ 
     \hline
     Working domain & Frequency &  Time and Frequency & Frequency \\ 
     \hline
     Neutron transport method & Diffusion & Diffusion and $\text{SP}_\text{N}$ & Diffusion\\ 
     \hline
     Energy group & Two group & Multigroup & Multigroup \\ 
     \hline
     Spatial discretization & Box-scheme finite difference (rectangular only) & Finite element & Finite volume (rectangular and triangular) \\ 
     \hline
     Code language & Matlab & C++ & Python\\ 
     \hline
    \end{tabular}} \quad
    \label{table:other_works}
\end{table}

In terms of the neutron noise unfolding method, \cite{demaziereIdentificationLocalizationAbsorbers2005} provided the theory and derivation of the noise unfolding method. The unfolding method is used to diagnose an absorber of variable strength noise source. This is further supported by \cite{hosseiniNoiseSourceReconstruction2014}, which provides several additional cases solved using a hybrid method. This work will focus on improving the existing method to solve the generalized noise source. Compared to both papers, the approach in this work focuses on neutron noise reconstruction, which can later be integrated with the known inverse of the Green’s function matrix.

\section{Outline}

This dissertation describes the motivation, objectives, methodology, results, and conclusions towards the advancement of the neutron noise unfolding methods and their application \gls*{HTGR} core.
The remainder of this dissertation is organized as follows.
Chapter \ref{ch:literature_review} presents a literature review of neutron noise methods and their applications for diagnostics.
Chapter \ref{ch:simulations} presents the development of computational tools for forward, adjoint, and noise calculations, as well as some application for verification purposes.
Chapter \ref{ch:noise_HTTR} focuses on the application of neutron noise methods for \gls*{AVS} and \gls*{FAV} in 2D and 3D \gls*{HTTR} core.
Chapter \ref{ch:unfolding_methods} introduces the novel methods of neutron noise unfolding and their application to 2D and 3D \gls*{HTTR} core. This includes the comparison between existing neutron noise unfolding methods and the novel methods.
Finally, Chapter \ref{ch:conclusions} summarizes the findings of this dissertation and provides recommendations for future work.
